% from ega2-1-fr.tex
\oldpage[II]{61}

On en conclut que lorsque $Y$ est un préschéma quelconque, les trois
conditions suivantes, pour un $\mathcal{O}_X$-Module quasi-cohérent
$\mathcal{F}$ tel que $p_*(\mathcal{F})$ soit \emph{quasi-cohérent}, sont
équivalentes~:
\begin{enumerate}
  \item[a)] \emph{L'homomorphisme canonique
    $\sigma:p^*(p_*(\mathcal{F}))\to\mathcal{F}$ est surjectif.}
  \item[b)] \emph{Il existe un $\mathcal{O}_Y$-Module quasi-cohérent
    $\mathcal{G}$ et un homomorphisme surjectif
    $p^*(\mathcal{G})\to\mathcal{F}$.}
  \item[c)] \emph{Pour tout ouvert affine $U$ de $Y$,
    $\mathcal{F}|p^{-1}(U)$ est engendré par ses sections au-dessus de
    $p^{-1}(U)$.}
\end{enumerate}
On vient en effet de prouver l'équivalence de \emph{a)} et \emph{c)}. D'autre
part, il est clair que \emph{a)} entraîne \emph{b)}, $p_*(\mathcal{F})$ étant
par hypothèse quasi-cohérent. Inversement, tout homomorphisme
$u:p^*(\mathcal{G})\to\mathcal{F}$ se factorise en
$p^*(\mathcal{G})\to p^*(p_*(\mathcal{F}))\xrightarrow{\sigma}\mathcal{F}$
(\hyperref[0.3.5.4.4-fr]{0, 3.5.4.4}), donc si $u$ est surjectif il en est de
même de $\sigma$, ce qui prouve que \emph{b)} entraîne \emph{a)}.

\begin{corollary}[3.4.8]
\label{II.3.4.8-fr}
Supposons vérifiées les hypothèses de (\hyperref[II.3.4.4-fr]{3.4.4}) et
supposons en outre que $Y$ soit un schéma quasi-compact ou que l'espace
sous-jacent à $Y$ soit noethérien. Soit $\mathcal{F}$ un
$\mathcal{O}_X$-Module quasi-cohérent de type fini~; il existe alors un entier
$n_0$ tel que pour $n\geq n_0$, $\mathcal{F}$ soit isomorphe à un quotient
d'un $\mathcal{O}_X$-Module de la forme $(p^*(\mathcal{G}))(-n)$ où
$\mathcal{G}$ est un $\mathcal{O}_Y$-Module quasi-cohérent de type fini
(dépendant de $n$).
\end{corollary}

\begin{proof}
Comme le morphisme structural $X\to Y$ est séparé et de type fini,
$p_*(\mathcal{F}(n))$ est quasi-cohérent
(\hyperref[I.9.2.2-fr]{I, 9.2.2, b)}), donc limite inductive de ses
sous-$\mathcal{O}_Y$-Modules quasi-cohérents de type fini, en vertu de
l'hypothèse sur $Y$ (\hyperref[I.9.4.9-fr]{I, 9.4.9}). On en déduit par
(\hyperref[II.3.4.6-fr]{3.4.6}),
(\hyperref[0.4.3.2-fr]{0, 4.3.2}) et
(\hyperref[0.5.2.3-fr]{0, 5.2.3}) que $\mathcal{F}(n)$ est l'image canonique
d'un $\mathcal{O}_X$-Module de la forme $p^*(\mathcal{G})$, où $\mathcal{G}$
est un sous-$\mathcal{O}_Y$-Module quasi-cohérent de type fini de
$p_*(\mathcal{F}(n))$~; le corollaire résulte alors de
(\hyperref[II.3.2.5.2-fr]{3.2.5.2}) et
(\hyperref[II.3.2.7.1-fr]{3.2.7.1}).
\end{proof}

\subsection{Comportements fonctoriels.}
\label{subsection:II.3.5-fr}

\begin{env}[3.5.1]
\label{II.3.5.1-fr}
Soient $Y$ un préschéma, $\mathcal{S}$, $\mathcal{S}'$ deux
$\mathcal{O}_Y$-Algèbres graduées quasi-cohérentes à degrés positifs~; posons
$X=\operatorname{Proj}(\mathcal{S})$, $X'=\operatorname{Proj}(\mathcal{S}')$,
et soient $p$, $p'$ les morphismes structuraux de $X$ et $X'$ dans $Y$. Soit
$\varphi:\mathcal{S}'\to\mathcal{S}$ un $\mathcal{O}_Y$-homomorphisme
d'Algèbres graduées. Pour tout ouvert affine $U$ de $Y$, posons
$S_U=\Gamma(U,\mathcal{S})$, $S'_U=\Gamma(U,\mathcal{S}')$~;
l'homomorphisme $\varphi$ définit un homomorphisme
$\varphi_U:S'_U\to S_U$ de $A_U$-algèbres graduées, en posant
$A_U=\Gamma(U,\mathcal{O}_Y)$. Il lui correspond dans $p^{-1}(U)$ un ensemble
ouvert $G(\varphi_U)$ et un morphisme
$\Phi_U:G(\varphi_U)\to p'^{-1}(U)$
(\hyperref[II.2.8.1-fr]{2.8.1}). En outre, si $V\subset U$ est un ouvert
affine, le diagramme
\[
\label{II.3.5.1.1-fr}
  \xymatrix{
    S'_U\ar[r]^{\varphi_U}\ar[d]
    & S_U\ar[d]\\
    S'_V\ar[r]_{\varphi_V}
    & S_V
  }
\tag{3.5.1.1}
\]
est commutatif, et on vérifie aussitôt, à partir des définitions
(\hyperref[II.2.8.1-fr]{2.8.1}), que l'on a
\[
  G(\varphi_V)=G(\varphi_U)\cap p^{-1}(V)
\]
et que $\Phi_V$ est la restriction de $\Phi_U$ à $G(\varphi_V)$. On a ainsi
défini une partie ouverte $G(\varphi)$ de $X$ telle que
$G(\varphi)\cap p^{-1}(U)=G(\varphi_U)$ pour tout ouvert affine $U\subset Y$,
et un $Y$-morphisme

